\chapter{Similaridade com o kernel Linux}
\label{cap:similaridade}

\section{O que significa ``similaridade''}

Comparar este projeto com o Linux exige separar tr\^es perguntas diferentes:

\begin{enumerate}
  \item \textbf{Similaridade funcional}: quais servi\c{c}os do Linux existem
        aqui e com que profundidade?
  \item \textbf{Similaridade arquitetural}: as solu\c{c}\~oes de boot,
        mem\'oria e drivers usam os mesmos conceitos e mecanismos de
        hardware?
  \item \textbf{Similaridade de c\'odigo}: houve c\'opia, tradu\c{c}\~ao
        linha a linha ou reutiliza\c{c}\~ao de algoritmos?
\end{enumerate}

As respostas s\~ao bastante diferentes. O projeto \'e conceitualmente
pr\'oximo do Linux no caminho de boot e no uso de VGA/PS/2, mas cobre uma
fra\c{c}\~ao min\'uscula do kernel completo. N\~ao h\'a c\'odigo copiado: a
implementa\c{c}\~ao \'e independente e escrita em Rust.

\section{Metodologia}

A compara\c{c}\~ao usa como refer\^encia a \'arvore Linux 7.3.1 presente em
\code{linux-7.3.1/}. Foram coletadas:

\begin{itemize}
  \item linhas f\'isicas dos 12 arquivos Linux mais pr\'oximos;
  \item linhas f\'isicas dos 19 m\'odulos Rust;
  \item funcionalidades presentes, ausentes e parciais;
  \item mecanismos de hardware usados (Multiboot2, GDT, PAE, long mode,
        VGA texto, 8042, UART 16550);
  \item presen\c{c}a ou aus\^encia de subsistemas maiores (IDT, alocador,
        scheduler, syscalls, FS, rede).
\end{itemize}

A pontua\c{c}\~ao \'e uma estimativa fundamentada, n\~ao uma m\'etrica formal.
Ela serve para comunicar ordem de grandeza e evitar tanto a conclus\~ao
ing\^enua de que ``\'e um Linux em Rust'' quanto a conclus\~ao oposta de que
``n\~ao tem nada a ver''.

\section{Correspond\^encia de arquivos}

\begin{longtable}{@{}p{7.2cm}r p{4.3cm}@{}}
\caption{Arquivos Linux an\'alogos e m\'odulos Rust correspondentes.}
\label{tab:correspondencia}\\
\toprule
\textbf{Arquivo Linux 7.3.1} & \textbf{Linhas} & \textbf{Equivalente Rust} \\
\midrule
\endfirsthead
\toprule
\textbf{Arquivo Linux 7.3.1} & \textbf{Linhas} & \textbf{Equivalente Rust} \\
\midrule
\endhead
\code{arch/x86/boot/header.S} & 619 & \code{boot.rs} (cabe\c{c}alho) \\
\code{arch/x86/boot/pm.c} & 125 & \code{boot.rs} (stub) \\
\code{arch/x86/kernel/head\_64.S} & 686 & \code{boot.rs} (modo longo) \\
\code{arch/x86/kernel/cpu/common.c} & 2.688 & \code{boot.rs} (GDT) + \code{arch.rs} \\
\code{arch/x86/mm/init\_64.c} & 1.557 & \code{boot.rs} (p\'aginas) \\
\code{drivers/video/console/vgacon.c} & 1.226 & \code{vga.rs} + \code{tty.rs} \\
\code{drivers/input/serio/i8042.c} & 1.664 & \code{keyboard.rs} \\
\code{drivers/input/keyboard/atkbd.c} & 1.981 & \code{keyboard.rs} \\
\code{drivers/tty/serial/8250/8250\_port.c} & 3.638 & \code{serial.rs} \\
\code{kernel/printk/printk.c} & 5.188 & macros + \code{tty.rs} \\
\code{kernel/panic.c} & 1.274 & \emph{panic handler} em \code{main.rs} \\
\code{init/main.c} & 1.706 & \code{kernel\_main} em \code{main.rs} \\
\code{mm/slab\_common.c}, \code{mm/slub.c} & --- & \code{heap.rs} \\
\code{fs/ramfs/inode.c} & --- & \code{ramfs.rs} \\
\code{fs/open.c}, \code{fs/file\_table.c} & --- & \code{vfs.rs} \\
\code{fs/proc/} & --- & \code{proc.rs} \\
\code{kernel/sched/core.c} & --- & \code{task.rs} \\
\code{arch/x86/entry/syscalls/syscall\_64.tbl} & --- & \code{syscall.rs} \\
\midrule
\textbf{Total Linux (12 primeiros arquivos)} & \textbf{22.352} & --- \\
\textbf{Total Rust (19 m\'odulos)} & \textbf{7.973} & --- \\
\bottomrule
\end{longtable}

As linhas marcadas com \code{---} ampliam o mapeamento conceitual depois da
adi\c{c}\~ao do heap, do VFS, do \code{/proc}, das tarefas e das syscalls de
arquivo; elas n\~ao fazem parte da compara\c{c}\~ao de linhas, que continua
usando os doze primeiros arquivos Linux.

A raz\~ao de linhas \'e:

\[
\frac{7.973}{22.352} \approx 35{,}67\%.
\]

Em rela\c{c}\~ao \`a \'arvore Linux completa, que tem aproximadamente
38.062.672 linhas de \code{.c}, \code{.h} e \code{.S} na vers\~ao usada, a
raz\~ao \'e:

\[
\frac{7.973}{38.062.672} \approx 0{,}0209\%.
\]

\begin{atencao}[Linhas de c\'odigo n\~ao medem similaridade]
Os 12 arquivos Linux n\~ao t\^em o mesmo escopo que os 19 m\'odulos Rust. O
\code{vgacon.c}, por exemplo, suporta v\'arias arquiteturas, consoles
virtuais e modos gr\'aficos; o \code{vga.rs} implementa apenas texto 80x25.
Al\'em disso, grande parte das linhas Linux \'e coment\'ario, configura\c{c}\~ao
e tratamento de casos de hardware. A compara\c{c}\~ao de LOC serve apenas como
ordem de grandeza.
\end{atencao}

\section{Similaridade funcional por dimens\~ao}

A tabela abaixo distribui 100 pontos entre dez dimens\~oes e atribui uma nota
de 0 a 100 para o quanto o projeto cobre da dimens\~ao correspondente no
Linux. A \'ultima coluna \'e a contribui\c{c}\~ao ponderada.

\begin{table}[H]
\centering
\caption{Pontua\c{c}\~ao de similaridade funcional e arquitetural.}
\label{tab:score}
\begin{tabular}{@{}lcrrr@{}}
\toprule
\textbf{Dimens\~ao} & \textbf{Peso} & \textbf{Linux} & \textbf{Projeto} & \textbf{Contrib.} \\
\midrule
Boot e protocolo de firmware & 10 & 100 & 30 & 3,00 \\
Transi\c{c}\~ao de modo e GDT & 8 & 100 & 35 & 2,80 \\
Pagina\c{c}\~ao e MMU & 12 & 100 & 25 & 3,00 \\
Console e v\'ideo & 15 & 100 & 25 & 3,75 \\
Teclado e entrada & 12 & 100 & 18 & 2,16 \\
Serial e tty & 5 & 100 & 20 & 1,00 \\
Interrup\c{c}\~oes e exce\c{c}\~oes & 12 & 100 & 30 & 3,60 \\
Mem\'oria f\'isica e heap & 8 & 100 & 25 & 2,00 \\
Processos, scheduler e syscalls & 12 & 100 & 45 & 5,40 \\
Arquivos, rede e outros drivers & 6 & 100 & 10 & 0,60 \\
\midrule
\textbf{Total} & \textbf{100} & --- & --- & \textbf{27,31} \\
\bottomrule
\end{tabular}
\end{table}

O \'indice ponderado \'e \textbf{27,3\%}. Ele deve ser lido como:

\begin{quote}
``Dentro das dez dimens\~oes escolhidas, o projeto reproduz cerca de 27,3\%
do que o Linux oferece, com peso maior para boot, console e teclado.''
\end{quote}

N\~ao significa que 27,3\% do c\'odigo do Linux foi reimplementado. Em
n\'umero de subsistemas e servi\c{c}os, a fra\c{c}\~ao \'e muito menor: o
projeto tem apenas um heap \emph{first-fit}, um escalonador preemptivo justo
e ponderado com \code{nice}, pesos, \code{vruntime}, \code{sleep} bloqueante,
filas de espera por evento, sinais no estilo Unix e colheita de zumbis, um
sistema de arquivos em mem\'oria e um \code{/proc} sint\'etico, enquanto o
Linux tamb\'em tem buddy allocator, slab, CFS/EEVDF completo com classes de
tempo real, uma infraestrutura de sinais e wait queues muito mais rica, um VFS
completo com block devices e rede.

\subsection{Leitura por dimens\~ao}

\begin{itemize}
  \item \textbf{Boot (30/100)}: o projeto usa Multiboot2 e GRUB, enquanto o
        Linux usa o protocolo de boot pr\'oprio, setup em tempo real,
        descompress\~ao e stubs EFI. O conceito de ``carregar um ELF e
        transferir o controle'' \'e o mesmo; a implementa\c{c}\~ao \'e bem
        mais simples.
  \item \textbf{Modo longo e GDT (35/100)}: a sequ\^encia PAE $\rightarrow$
        LME $\rightarrow$ PG \'e a mesma do \code{head\_64.S}; h\'a GDT com
        ring 3 e TSS, mas faltam \emph{per-CPU}, \emph{stack canary},
        \emph{paravirt} e tratamento de CPUs secund\'arias.
  \item \textbf{Pagina\c{c}\~ao (25/100)}: h\'a \emph{huge pages} de 2 MiB no
        boot, alocador de frames por bitmap, tabelas de 4 KiB e demand
        paging; faltam buddy allocator, \code{vmalloc}, \code{mmap}, COW,
        swap, NUMA e prote\c{c}\~ao \code{NX}.
  \item \textbf{Console (25/100)}: o modo texto VGA e a rolagem s\~ao
        an\'alogos ao \code{vgacon} e h\'a oito consoles virtuais com troca
        por \code{Alt+F1}..\code{Alt+F8}; faltam fontes, Unicode,
        \emph{fbcon}, \code{printk} com n\'iveis e buffer de log.
  \item \textbf{Teclado (18/100)}: o acesso ao 8042 e a tradu\c{c}\~ao de
        scancodes s\~ao semelhantes, mas faltam IRQ1, \emph{input core},
        \code{evdev}, LEDs, layouts e teclas estendidas.
  \item \textbf{Serial (20/100)}: a programa\c{c}\~ao b\'asica da 16550 \'e
        parecida, mas faltam \emph{tty layer}, \emph{termios}, IRQs, FIFO
        de entrada, DMA e controle de fluxo.
  \item \textbf{Interrup\c{c}\~oes (30/100)}: h\'a IDT, \emph{stubs} para
        exce\c{c}\~oes, PIC remapeado e PIT em IRQ0; faltam APIC, IRQ1,
        IRQs esp\'urias, \emph{softirq}, workqueues e preemp\c{c}\~ao.
  \item \textbf{Mem\'oria/heap (25/100)}: h\'a alocador de frames por bitmap,
        demand paging e um heap \emph{first-fit} de 1 MiB atr\'as de
        \code{\#[global\_allocator]}, mas n\~ao h\'a buddy allocator, slab,
        page cache, \code{vmalloc}, swap ou NUMA.
  \item \textbf{Processos (45/100)}: h\'a tabela de tarefas com PID/PPID e
        estados, troca de contexto que preserva \code{RFLAGS}, escalonador
        preemptivo justo e ponderado com \code{nice}, pesos no estilo
        \code{sched\_prio\_to\_weight} e \code{vruntime}, \code{sleep}
        bloqueante com sem\^antica de \code{schedule\_timeout}, filas de
        espera por evento com \code{wait\_event}/\code{wake\_up}, sinais no
        estilo Unix com m\'ascaras de pendentes/bloqueados, disposi\c{c}\~oes e
        a\c{c}\~oes padr\~ao, tarefas paradas/continuadas, status de sa\'ida,
        zumbis, \code{wait}/\code{SIGCHLD} e doze syscalls; n\~ao h\'a
        \code{fork}/\code{exec}, espa\c{c}os de endere\c{c}o separados,
        classes de tempo real nem \'arvore rubro-negra completa do CFS/EEVDF.
  \item \textbf{FS/rede (10/100)}: h\'a ramfs, VFS com dezesseis descritores e
        um \code{/proc} sint\'etico, mas n\~ao h\'a block layer, driver de
        disco, tabela de mounts, modelo de permiss\~oes, page cache, TCP/IP ou
        netfilter.
\end{itemize}

\section{Similaridade arquitetural}

A tabela a seguir compara mecanismos concretos, independentemente do tamanho
do c\'odigo.

\begin{table}[H]
\centering
\caption{Compara\c{c}\~ao de mecanismos de hardware e arquitetura.}
\label{tab:mecanismos}
\begin{tabularx}{\linewidth}{@{}p{3.6cm}p{4.7cm}X@{}}
\toprule
\textbf{Mecanismo} & \textbf{Linux 7.3.1} & \textbf{Rust VGA Kernel} \\
\midrule
Carregamento & Protocolo Linux + setup + descompress\~ao & Multiboot2 via GRUB \\
Entrada inicial & \code{startup\_32}/\code{startup\_64} & \code{\_start} em assembly \\
Modo longo & PAE + LME + \code{CR0.PG} & Mesma sequ\^encia, simplificada \\
GDT & \code{gdt\_table}, \emph{per-CPU}, TSS & 7 descritores, ring 3 e
TSS/IST \\
Pagina\c{c}\~ao & 4/5 n\'iveis, alocadores, \code{vmalloc} & 4 n\'iveis,
2 MiB \emph{huge pages} + folhas 4 KiB sob demanda \\
Console & \code{vgacon}, VT, \code{fbcon}, \code{printk} & \code{vga.rs}
renderizador + \code{tty.rs} com oito consoles \\
Teclado & i8042 + serio + atkbd + input & 8042 por polling (IRQ1 mascarada) \\
Serial & 8250 + tty + termios & 16550 por polling \\
Entrada & IRQ1/IRQ4, \emph{softirq}, workqueues & IRQ0 (timer) + polling no
shell + \code{int 0x80} \\
Tempo & PIT/APIC, \code{jiffies}, timers & PIT a ~100 Hz, contador de
\emph{ticks}, contabilidade ponderada de \code{vruntime} \\
Mem\'oria & buddy, slab, \code{vmalloc}, page cache & Bitmap de frames,
demand paging, heap \emph{first-fit} \\
Processos & CFS/EEVDF, namespaces, cgroups & Tabela de tarefas,
escalonador justo e ponderado (\code{nice}/peso/\code{vruntime}), \code{sleep}
bloqueante, wait queues, sinais e \code{wait}/zumbis; sem \code{fork} ou
namespaces \\
Syscalls & \code{syscall}/\code{sysret}, \code{ptrace} & \code{int 0x80} com
doze chamadas (\code{exit}..\code{kill}) \\
Arquivos & VFS, inodes, dentries, page cache & ramfs, 16 descritores,
\code{/proc} sint\'etico; sem block layer \\
Rede & sockets, TCP/IP, netfilter & Ausente \\
\bottomrule
\end{tabularx}
\end{table}

\section{Similaridade de c\'odigo}

\begin{itemize}
  \item \textbf{C\'opia de c\'odigo}: nenhuma. O projeto foi escrito em Rust
        e n\~ao cont\'em trechos traduzidos do C ou do assembly do Linux.
  \item \textbf{Algoritmos}: as sequ\^encias de inicializa\c{c}\~ao do 8042,
        UART e modo longo seguem a documenta\c{c}\~ao de hardware e padr\~oes
        consolidados, que o Linux tamb\'em segue. Isso \'e converg\^encia
        t\'ecnica, n\~ao c\'opia.
  \item \textbf{Constantes}: valores como \code{0x3F8}, \code{0x60},
        \code{0x64}, \code{0x3D4}, \code{0xB8000} e os bits de \code{CR0},
        \code{CR4} e \code{EFER} s\~ao definidos pelo hardware e aparecem em
        qualquer driver, inclusive no Linux.
  \item \textbf{Licen\c{c}a}: o projeto se declara MIT/Apache-2.0. Se
        trechos GPL-2.0 do Linux fossem copiados, a licen\c{c}a precisaria
        ser compat\'ivel com GPL. Como a implementa\c{c}\~ao \'e
        independente, n\~ao h\'a essa obriga\c{c}\~ao.
\end{itemize}

\section{Similaridade global}

\begin{table}[H]
\centering
\caption{Resumo das diferentes medidas de similaridade.}
\label{tab:resumo-similaridade}
\begin{tabularx}{\linewidth}{@{}lX@{}}
\toprule
\textbf{Medida} & \textbf{Resultado} \\
\midrule
C\'odigo copiado & 0\% (implementa\c{c}\~ao independente) \\
Linhas vs.\ 12 arquivos an\'alogos & $\approx 35{,}67\%$ \\
Linhas vs.\ \'arvore Linux completa & $\approx 0{,}0209\%$ \\
Similaridade funcional ponderada (10 dimens\~oes) & $\approx 27{,}3\%$ \\
Cobertura de subsistemas do kernel completo & muito baixa; ordem de 2\% \\
Similaridade conceitual no boot & alta (mesmos mecanismos x86) \\
Similaridade conceitual em drivers & m\'edia/baixa (subconjunto por polling) \\
Similaridade conceitual em servi\c{c}os & baixa (MM, scheduler e VFS b\'asicos; sem rede) \\
\bottomrule
\end{tabularx}
\end{table}

\subsection{Classifica\c{c}\~ao}

Em uma escala qualitativa:

\begin{itemize}
  \item \textbf{Relacionamento com o Linux}: ``inspirado nos mesmos
        conceitos e no mesmo hardware'', n\~ao ``porte'' nem ``clone''.
  \item \textbf{Similaridade global}: \textbf{baixa}.
  \item \textbf{Valor did\'atico}: alto para entender boot, modo longo,
        VGA, PS/2 e serial; baixo para entender o Linux como sistema
        operacional completo.
  \item \textbf{Nota de 0 a 10 para ``quanto do Linux foi reimplementado''}:
        aproximadamente \textbf{2,7/10} na pontua\c{c}\~ao ponderada; menos de
        \textbf{0,1/10} se o crit\'erio for o kernel completo em
        funcionalidades.
\end{itemize}

\begin{nota}[Como usar esta an\'alise]
Se o objetivo for estudar o Linux, este projeto serve como ``mapa do
caminho inicial'': ele mostra, em c\'odigo pequeno, o que acontece antes de
o kernel Linux come\c{c}ar a usar seus subsistemas. Se o objetivo for
substituir o Linux, o projeto est\'a a v\'arias ordens de grandeza de
dist\^ancia: faltam CFS/EEVDF, \code{fork}/\code{exec}, espa\c{c}os de
endere\c{c}o por processo, sistema de arquivos em disco, rede, drivers e
endurecimento de seguran\c{c}a.
\end{nota}

\section{Diferen\c{c}as essenciais}

\begin{table}[H]
\centering
\caption{Diferen\c{c}as essenciais entre o projeto e o Linux.}
\label{tab:diferencas}
\begin{tabularx}{\linewidth}{@{}p{3.2cm}XX@{}}
\toprule
\textbf{Aspecto} & \textbf{Linux} & \textbf{Rust VGA Kernel} \\
\midrule
Linguagem & C e assembly & Rust e assembly \\
Tamanho & $\sim$38 milh\~oes de linhas & 7.973 linhas \\
Modelo de boot & Protocolo pr\'oprio, setup, EFI & Multiboot2/GRUB \\
Interrup\c{c}\~oes & IDT, PIC/APIC, timers & IDT, PIC/PIT; teclado em polling \\
Concorr\^encia & SMP, preemp\c{c}\~ao, RCU & Uma CPU, escalonador preemptivo
justo e ponderado com sleep bloqueante e wait queues \\
Mem\'oria & Buddy, slab, VM, page cache & Bitmap de frames, demand paging,
heap \emph{first-fit} \\
Processos & Scheduler, sinais, namespaces, cgroups & Tabela de tarefas,
escalonador justo e ponderado, wait queues, sinais e \code{wait}/zumbis; sem
namespaces \\
Entrada/sa\'ida & VFS, block, net, tty & ramfs, VFS, \code{/proc}, portas de I/O diretas \\
Seguran\c{c}a & Rings, SMEP/SMAP, KASLR, LSM & Rings 0/3 e U/S; sem NX,
SMEP ou SMAP \\
Portabilidade & V\'arias arquiteturas & Apenas x86\_64 BIOS \\
Licen\c{c}a & GPL-2.0 & MIT/Apache-2.0 \\
\bottomrule
\end{tabularx}
\end{table}
